% ============================================================
%  Section 8 --- Difficultés rencontrées et corrections
% ============================================================

\section{Difficultés rencontrées et évolutions méthodologiques}
\label{sec:difficulties}

Ce stage a été marqué par plusieurs difficultés techniques et méthodologiques, dont
certaines n'ont été identifiées qu'après des semaines d'expérimentation. Cette section
les documente honnêtement, en expliquant comment chacune a été détectée, diagnostiquée
et résolue --- ou pourquoi elle reste une limitation ouverte.

\subsection{Fuite d'information inter-splits (correction H1/H3)}
\label{sec:diff_leakage}

\paragraph{Problème.}
La version initiale du code d'évaluation appliquait \texttt{fit\_predict} indépendamment
sur chaque split pour assigner les labels de cluster. En pratique, le clustering
sur le split de validation et le clustering sur le split de test produisaient chacun leur
propre géométrie --- sans lien avec les centroides calculés sur train.
Cela donnait des scores de silhouette artificiellement optimistes
(\texttt{sil\_val}~$= 0.601$, \texttt{sil\_test}~$= 0.615$) car le clustering était
ajusté à la distribution locale de chaque split plutôt qu'évalué en généralisation.

\paragraph{Détection.}
L'incohérence a été détectée lors de l'analyse des clusters val et test : les centroides
val ne correspondaient pas aux centroides train, et des épisodes de même scénario
se retrouvaient dans des clusters différents selon le split. Ce comportement est un
signal clair de fuite d'information.

\paragraph{Correction.}
La correction consiste à calculer les centroides une seule fois sur train, puis à
assigner les labels val et test par \textbf{plus proche centroïde} (\emph{nearest centroid})
depuis ces centroides figés. Cette approche simule fidèlement le comportement en
production : un nouvel épisode est typé par rapport à la structure apprise sur les
données d'entraînement, sans re-clustering.
Après correction, les scores descendent à \texttt{sil\_test}~$= 0.414$ (graine 42),
soit une baisse de $-0.201$. Malgré cette baisse, H1 reste validée --- ce qui renforce
la confiance dans la robustesse réelle du résultat.

La même correction a été appliquée à H3 : la sélection de $k^*_i$ utilisait
précédemment le jeu de test, ce qui constituait une forme d'optimisme par sélection.
$k^*_i$ est désormais sélectionné exclusivement sur val, et l'AUROC rapporté
sur test uniquement.

\subsection{Bug LayerNorm TCN (correction silencieuse)}
\label{sec:diff_layernorm}

\paragraph{Problème.}
Un commit de refactoring (commit \texttt{6543c69}) avait activé un module
\texttt{LayerNorm} dans le TCN de l'encodeur après l'entraînement : le modèle était
entraîné sans LayerNorm activé, mais l'inférence l'utilisait.
Ce décalage train/inférence produisait des embeddings légèrement hors distribution,
dégradant silencieusement les performances sans déclencher d'erreur.

\paragraph{Détection.}
Le bug a été identifié lors de la comparaison des résultats avant et après le commit :
l'AUROC H3 avait baissé de $1$--$2$\,pp sans modification du protocole d'évaluation.
Une inspection du code a révélé la divergence entre le forward pass à l'entraînement
et à l'inférence.

\paragraph{Correction.}
Le forward pass de l'encodeur a été restauré au comportement de l'entraînement
(LayerNorm désactivé). Les précurseurs ont été réentraînés à partir de l'encodeur corrigé.
L'AUROC H3 a retrouvé ses niveaux attendus (+0.3 à +1.6\,pp selon le cluster).
Ce bug illustre l'importance de versionner séparément les checkpoints du modèle
et le code d'inférence, ou d'utiliser \texttt{torch.jit.script} pour figer le graphe
de calcul au moment de l'entraînement.

\subsection{Méthode d'explicabilité gradient invalidée}
\label{sec:diff_saliency}

\paragraph{Problème.}
La méthode d'explicabilité initiale pour les fiches de clusters était le gradient saliency
(gradient$\times$input) : pour chaque feature, le gradient du score de clustering
par rapport à cette feature est calculé, puis multiplié par sa valeur.
Cette méthode est rapide et différentiable, ce qui la rendait attractive pour générer
automatiquement des fiches d'interprétabilité.

\paragraph{Détection.}
La validation croisée entre gradient$\times$input et permutation importance a révélé
une corrélation de Spearman $\rho = -0.34$ : les deux méthodes s'accordent sur quelles
features sont importantes dans seulement $33$\,\% des cas, et sont \emph{anti-corrélées}
en moyenne. Ce résultat était inattendu et a nécessité une investigation approfondie.

\paragraph{Cause identifiée.}
Le problème vient de la non-linéarité de l'encodeur STGCN combinée à la perte
contrastive siamoise. Le gradient mesure la sensibilité locale du score en un point,
mais cette sensibilité locale ne reflète pas l'importance globale d'une feature pour
la structuration des clusters --- d'autant plus que la loss contrastive crée des
gradients qui pointent dans des directions géométriquement complexes dans l'espace
des embeddings.

\paragraph{Correction.}
Les fiches ont été régénérées avec la \textbf{permutation importance} (50 shuffles) :
chaque feature est permutée aléatoirement et la baisse de score de silhouette mesure
son importance. Cette méthode est model-agnostic et mesure une importance globale,
non locale. Elle a ensuite été validée par KernelSHAP \cite{lundberg2017}
($\rho_\text{Spearman}(\text{SHAP}, \text{permutation}) > 0$ pour $9/10$ clusters),
ce qui confirme la cohérence des fiches pour la majorité des types.

\subsection{Biais écologique dans la Transfer Entropy cluster-level}
\label{sec:diff_te}

\paragraph{Problème.}
La première implémentation de la Transfer Entropy pour l'ontologie construisait une
seule série temporelle par cluster en concaténant tous les épisodes du cluster,
puis estimait la TE entre services sur cette série unique.
Cette approche confond deux niveaux d'analyse : les dynamiques intra-épisode
(comment les services interagissent pendant un incident) et les dynamiques inter-épisodes
(tendances à travers les répétitions d'un scénario).
C'est un biais écologique classique en écologie statistique, ici transposé au domaine
des séries temporelles.

\paragraph{Détection et correction.}
Le biais a été identifié lors de la comparaison des résultats entre un dry-run
à 20 permutations (qui produisait 2 relations causales, C6$\to$C8 et C2$\to$C8)
et 100 permutations (qui produisait 0 relation) : les relations observées à faibles
permutations étaient des faux positifs créés par la structure artificielle de la
concaténation.

La correction a été de basculer vers un estimateur \textbf{service-level} :
la TE est estimée indépendamment sur chaque épisode, puis les p-values sont agrégées
par correction BH-FDR et les TE moyennées par cluster.
Cette approche respecte l'indépendance des épisodes et produit 124 relations causales
significatives sur 8/10 clusters, dont le résultat validant que les drifts bénins
(C5, C6) ne génèrent aucune cascade causale inter-services.

\subsection{Résultats inattendus sur les baselines}
\label{sec:diff_baselines}

\paragraph{M\_only bat le modèle full pour H1.}
L'ablation par réentraînement a produit un résultat contre-intuitif : le modèle
entraîné uniquement sur les métriques Prometheus ($n_\text{feat} = 7$) surpasse le
modèle complet pour la silhouette de clustering ($+0.058$). Intuitivement, davantage
de features devraient aider.

L'explication rétrospective est un sur-paramétrage du typage contrastif sur
$n_\text{train} = 209$ épisodes : les features de traces et de logs introduisent
de la variabilité intra-cluster (signatures variables d'un épisode à l'autre pour un
même type de panne) qui brouille la perte contrastive. Sur ewat\_v4 avec
$n_\text{train}$ plus grand, ce comportement devrait s'inverser.
La leçon opérationnelle est de ne pas supposer que davantage de features améliore
nécessairement le clustering sur des petits datasets.

\paragraph{STGCN n'améliore pas l'AUROC macro sur les scénarios.}
La comparaison B3/B4 (features brutes vs STGCN sur les scénarios Chaos Mesh)
a montré un macro-AUROC identique ($0.835$). Cela a d'abord semblé indiquer que
l'encodeur n'apportait rien. Une analyse détaillée par scénario a révélé une
redistribution non triviale : $+0.270$ sur \texttt{fail\_slow\_cpu},
$-0.246$ sur \texttt{noisy\_neighbor}. La coïncidence numérique résulte d'une
compensation exacte sur $n = 45$ épisodes, et non d'une neutralité de l'encodeur.
Cette distinction entre neutralité et redistribution compensée est importante pour
interpréter les résultats sur ewat\_v4, où la compensation ne tiendra probablement plus.

\paragraph{Look-through moins performant que le seuil simple pour le TPR drift.}
H2a a montré que le look-through réduit le TPR drift de $0.67$ à $0.42$.
Ce résultat paraît paradoxal : le mécanisme devrait améliorer la détection, pas la
dégrader. L'explication est que le warm-up du DriftDetector ($\approx 8$ steps) et la
confirmation temporelle ($3$--$6$ steps) retardent le flag drift au point de le manquer
sur des épisodes courts --- là où le seuil simple déclenche immédiatement.
Le look-through est conçu pour réduire les faux positifs, non pour améliorer le TPR
sur des épisodes aussi courts.

\subsection{Collecte de données : pannes d'infrastructure}
\label{sec:diff_infra}

\paragraph{Nœud NotReady.}
Le nœud \texttt{observit-cluster1-workers-58w74-mwxb2} est resté en état NotReady
pendant toute la durée de la collecte ewat\_v3, rendant les métriques \texttt{disk\_io}
du service \texttt{product-catalog} indisponibles ($16.7$\,\% de NaN).
Ce problème n'a été identifié qu'après les premières validations du dataset, lorsque
le taux de NaN anormalement élevé pour une seule combinaison service/feature a attiré
l'attention.

L'imputation par zéro a été retenue comme solution de compromis : elle n'introduit pas
de biais directionnel (zéro est une valeur neutre pour les IOPS), mais sous-estime
l'importance de la feature --- confirmé par l'ablation LOO ($\Delta = -0.088$).
La correction structurelle (remplacement du nœud ou déploiement OTel SDK applicatif)
est planifiée pour ewat\_v4.

\paragraph{Épisode \texttt{network\_loss\_018} rejeté.}
Un épisode a produit $100$\,\% de NaN sur toutes les features de logs (Loki indisponible
pendant la fenêtre d'enregistrement). Cet épisode a été rejeté, réduisant le dataset
de 300 à 299 épisodes. Ce type d'échec est inhérent à la collecte sur infrastructure
réelle : la stack d'observabilité elle-même peut être affectée par les injections
de pannes réseau.

\subsection{Transfert de domaine : hypothèse initiale invalidée}
\label{sec:diff_transfer}

L'hypothèse de départ pour l'évaluation sur RCAEval était qu'un simple ajustement
du scaler suffirait à transférer le pipeline sur un nouveau cluster hébergeant le même
logiciel. Cette hypothèse s'est révélée incorrecte : H3 AUROC reste à $\approx 0.5$
quel que soit $n_\text{few}$ avec la Stratégie A.

Le diagnostic a nécessité trois étapes. D'abord, tester quatre stratégies de
normalisation pour isoler l'effet du scaler. Ensuite, vérifier que l'encodeur
produisait bien des embeddings cohérents (H1 PASS avec instance norm + M\_only).
Enfin, vérifier que l'échec de H3 n'était pas dû à une inadéquation des features mais
à la durée des épisodes ($\times 2.3$) qui déplace les signatures pré-injection
hors de la fenêtre $k^*$ apprise sur ewat\_v3.

Ce diagnostic multi-factoriel a pris plusieurs jours et illustre la complexité du
domain shift dans les pipelines de machine learning sur séries temporelles :
le shift n'affecte pas uniformément toutes les étapes du pipeline.
